% page 209 du fac-simile (image p-208 du scan)
% bloc 172.7 x 250.6 mm ; marge gauche 29.3 mm
\pgc{7.620mm}{0.000mm}{
\l{}
\l{}
\l{}
\l{\cel{55}2Ol}
\l{}
\l{\sou{gratular}. (\sou{trans., pri, pro}). Dicar politeso-vorti (ad ulu) pri}
\l{\cel{3}la fortuno quan lu profitas. - DEFIS.}
\l{}
\l{\sou{grava}.\cel{2}I.\cel{2}Desfacile movebla, levebla, pro lua pezo. - II. (\sou{muz.})}
\l{\cel{3}Qua apartenas a la gradi infra di la muziko-skalo. - III.}
\l{\cel{3}(\sou{gram}.)En ula lingui (kom ex. la Franca)Qualeso di la supersig-}
\l{\cel{3}no analoga a komo inversigita qua donas a la vokalo \textquotedbl{}e\textquotedbl{} son-uno}
\l{\cel{3}apertita, o distingas ula homonimi (kom ex.en F.: la\cel{1}\sou{e là)}\cel{1}.}
\l{\cel{3}IV. (Afero) ne-vana, nefrivola, egardenda kom importoza. -EFIS.}
\l{}
\l{\sou{gravelo}. (\sou{patol.})\cel{2}Konkreciono qua naskas en la reni. - EF.}
\l{}
\l{\sou{gravio}. I. Sablo grosa-grana, qua devenas de la desagrego di roki}
\l{\cel{3}stonoza. - II. Sabl-uni grosa per qui onu kovras la alei di}
\l{\cel{3}gardeno. - EFR.}
\l{}
\l{\sou{gravida}. Qua portas filio (yuno) en sua ventro. - IL.}
\l{}
\l{\sou{gravitar}.\cel{2}(\sou{netrans.}) Obediar la forco pro qua la molekuli inter-}
\l{\cel{3}atraktas proporcionale ye lia maso. - DEFIS.}
\l{}
\l{\sou{greftar}. (\sou{trans.,}\cel{1}sur) Insertar sproso di planto sur altra planto,}
\l{\cel{3}por ke ica portez la floro, la frukto diita. - EF.}
\l{}
\l{\sou{grelar}.\cel{2}(\sou{netrans.})\cel{2}Falo di pluvo-guteti konjelita, grano-forma,}
\l{\cel{3}partikulare dum sturmo-vetero. - F.}
\l{}
\l{\sou{gremio}. I. Che persono sidanta : parto di lua korpo qua iras de}
\l{\cel{3}la tayo anla genui. - II. (\sou{metaf}.) La sino. - DIS.}
\l{}
\l{\sou{grenado}. I. (\sou{bot.)}\cel{2}Frukto sferoida qua kontenas semini reda, enklo-}
\l{\cel{3}zita en mikra celuli, produktita da arbusto de la familio}
\l{\cel{3}\textquotedbl{}mirtacei\textquotedbl{}. L.\cel{1}\sou{punica granatum.}\cel{2}- II. (\sou{militarto)}. Projektilo}
\l{\cel{3}grenadatra, bulo de fero, plena ye stupo e de pulvero, di qua onu}
\l{\cel{3}acendis la mecho por lansar lu manue. - DEFIRS.}
\l{}
\l{\sou{greso}. (\sou{mineral.}) Roko ek sabl-uni quarcozz. - II. Rekuro-sablo,}
\l{\cel{3}facita ek ta petro pulverigita. - FI.}
\l{}
\l{\sou{greto}. Asemblajo ajuroza ek stangi de fero o de ligno, uzata}
\l{\cel{3}kom klozilo di gardeno, parko, e c. - EFS.}
\l{}
\l{\sou{gribo}. Peco de panikulo o di lardo de porko, fritita en padelo.}
\l{}
\l{\sou{grifono}. I. (\sou{mitol}.) Animalo mi-agla, mi-leona. - II. (\sou{zool.})}
\l{\cel{3}Mikra hundo pudala, kun longa pili herisita sur la kapo e la}
\l{\cel{3}parto avana di la korpo. - DEFIS.}
\l{}
\l{\sou{grilar}. (\sou{trans}.) I. (\sou{koqua}rto). Bruletar, per fairo vivaca, nutrivo}
\l{\cel{3}(karno, fisho, pano) pozita sur koquo-utensilo ek dina stangi}
\l{\cel{3}de fero. - II. Submisar erco a fairo por igar lu plu friabla.}
\l{\cel{3}III. Submisar a la efiko di flamo kotonajo (pos texto) o}
\l{\cel{3}pultruno por supresar la lanugo. - EF.}
}
